\section{Conclusion}
In our work, we have introduced a new architecture that extends Algorithm Distillation (AD) for environments with variable action spaces, achieving invariance to their structure and size. Our approach consists of discarding the last linear layer, granting invariance to the action space structure, and making the model infer action semantics from the context, preparing it for the introduction of novel actions. We demonstrated Headless-AD's capability to generalize across new action spaces on a set of environments. We also observed its performance gains over vanilla AD, especially on larger action spaces. We hypothesize that this is due to the augmentation of the dataset by random embeddings, which was shown to improve the generalization abilities of agents \cite{kirsch2023towards}. Headless-AD marks the progress toward versatile foundational models in RL, ones capable of operating across an expanded range of environments. We hope that Headless-AD inspires further development of models that can adapt to any action space beyond discrete ones.

\textbf{Limitations.} Headless-AD shares AD's limitation of fixed sequence lengths, which may limit its effectiveness in environments with long episodes. A unique constraint of Headless-AD is its limit on the number of actions, dictated by the dimensionality of the action embeddings -- only as many actions as there are dimensions can be orthogonally represented. Going over this limit causes embeddings to become linearly dependent, unintentionally distributing probability across multiple actions. While action spaces in RL environments typically do not become excessively large, and increasing the dimension of embeddings could mitigate this issue, it remains a point for consideration.

\textbf{Future Work.} In our study, we demonstrated the ability of our algorithm to generalize to new action spaces, as shown by its performance on elementary tasks. To extend and validate these findings, future research should focus on more complex environments. This will offer a deeper insight into the algorithm's versatility and robustness in diverse and complex settings.

Additionally, to ensure the broader applicability and adaptability of our approach, it is essential to examine its compatibility and performance with various models beyond Algorithm Distillation (AD) \citep{laskin2022context}, such as the Decision Pretrained Transformer (DPT) \citep{lee2023supervised}. This exploration will provide a more comprehensive understanding of the algorithm's strengths and limitations, potentially leading to further improvements and a wider scope of applications in different contexts. 